\documentclass[a4paper,12pt,fleqn,headings=small,listof=totoc]{scrreprt}  %fleqn setzt die Formeln linksbündig

\usepackage[T1]{fontenc}
\usepackage[utf8]{inputenc}
\usepackage[ngerman]{babel}
\usepackage{helvet,courier,times}  % mit "times" statt "mathptmx" besserer Formelstil
\usepackage{amssymb} % für mathematische Symbole
\usepackage{stmaryrd} % beinhaltet das Widerspruchszeichen Blitz: \lightning
\usepackage{amsmath} % Mathematik
\usepackage{wasysym} % wasysym Zeichen darstellbar
\usepackage{eurosym} % Eurozeichen
\usepackage{enumitem}
\usepackage{array} 
\usepackage{icomma}  % no space after comma in math-mode, but space in vectors
\usepackage{xcolor, graphicx}
\usepackage{mathcomp}
\usepackage{cancel}
\usepackage{nicefrac} % Bruch als 3/4 darstellbar
\usepackage{comma} % Komma wird erkannt
\usepackage{hhline} % hhline Paket für Tabellen
\usepackage{booktabs}
\usepackage{pdfpages}			
\usepackage{blindtext}
\usepackage{verbatim}
\usepackage{graphicx}
\usepackage{multirow}
\usepackage{natbib}

\DeclareOldFontCommand{\bf}{\normalfont\bfseries}{\mathbf}					
						
\parindent 0pt %kein Einrücken des Absatzes

\usepackage[left=2cm,right=4cm,bottom=3cm,top=3cm]{geometry}  % Einstellungen der Seitenränder % Seitenbreite 

\usepackage{setspace}
\onehalfspacing  % für Zeilenabstand 1,5					
						
%Abstand zwischen Formeln verkleinern
\makeatletter
\g@addto@macro\normalsize{%
	\setlength{\abovedisplayskip}{0pt}
	\setlength{\belowdisplayskip}{8pt}
	\setlength{\abovedisplayshortskip}{0pt}
	\setlength{\belowdisplayshortskip}{0pt}
	}
\makeatother
									
%keine neuen Seiten im Literaturverzeichnis
\makeatletter          
\renewcommand*\bib@heading{%
	\section*{\bibname}%
	}
\makeatother
										
%ändert den Abstand von der Kapitelüberschrift zum oberen Blattrand
\renewcommand*\chapterheadstartvskip{\vspace{-1cm}}
						

%%%%%%%%%%%%%%%%%%%%%%%%%%%%%%%%%%%%%%%%%%%%%%%%%%%%%%%%%%%%%%%%%%%%%%%%%%%%%%%%%%%%%%%%%%%%%%


\begin{document}
\definecolor{MyDarkBlue}{rgb}{0.2,0,0.55}
\titlepage % Titelseite
\sffamily

\begin{center}
\includegraphics[scale=0.06]{feu-logo.eps}
\huge{\color{MyDarkBlue}{{FernUniversität in Hagen}}}
\end{center}

\vspace{1cm}

\begin{center}\Large
\textsc{Fakultät für Wirtschaftswissenschaft}
\end{center}

\vspace{1cm}

\begin{center}\Large
\textsc{Verfahren zur Klassifikation multivariater~Daten}
\end{center}

\vspace{2cm}

\begin{center}

 Diplomarbeit 
 
 Lehrstuhl für Betriebswirtschaftslehre,
  
 insbesondere Quantitative Methoden und 
 Wirtschaftsmathematik 

\end{center}

						
\begin{center}
\includegraphics[scale=0.8]{quam-logo.eps}
\end{center}

\vspace{4cm}

\begin{tabular}{lp{3.5em}l}
					&&\\
					&&\\
Vorname, Name: && Karsten Opitz\\
Matrikelnummer: && 6811221 \\
\end{tabular}  

% Titelseite Ende

\rmfamily % Schriftart
\newpage % neue Seite
\normalsize % Schriftgröße
\pagenumbering{roman} % Seitenzahlen mit kleinen römischen Ziffern

\setcounter{tocdepth}{3}
\tableofcontents % Inhaltsverzeichnis
\newpage % neue Seite


\listoffigures % Abbildungsverzeichnis
%\newpage % neue Seite

\listoftables % Tabellenverzeichnis
\newpage % neue Seite

\pagenumbering{arabic} % Seitenzahlen mit arabischen Ziffern


\chapter{Einführung}


\chapter{Klassifikationsverfahren im Data Mining} 
\label{chap:classification_methods}

\section{Data Mining}

\section{Klassifikationsverfahren}

\subsection{Statistische Verfahren}

\subsection{Decision Trees}

\subsection{Neurale Netze und Deep Learning}

\chapter{Entscheidungsunterstützung durch Klassifikationsverfahren}

\chapter{Werkzeuge}
\section{R}
\section{Python}
\section{KNIME}
\section{Weka}


\chapter{Schlussfolgerungen}
\label{chap:conclusions}
%%%%%%%%%%%%%%%%%%%%%%%%%% Anhang %%%%%%%%%%%%%%%%%%%%%%%%%

\nocite{*}
\bibliography{literatur}
\bibliographystyle{quam-natdin}	

\addcontentsline{toc}{chapter}{Literaturverzeichnis} % Literaturverzeichnis wird im Inhaltsverzeichnis aufgenommen

\begin{comment}

\newpage % neue Seite


\chapter*{Erklärung}
\addcontentsline{toc}{chapter}{Erklärung} % Eidesstattliche versicherung wird im Inhaltsverzeichnis aufgenommen

Ich erkläre, dass ich die Seminararbeit selbstständig und ohne unzulässige Inanspruchnahme Dritter verfasst habe. Ich habe dabei nur die angegebenen Quellen und Hilfsmittel verwendet und die aus diesen wörtlich, inhaltlich oder sinngemäß entnommenen Stellen als solche den wissenschaftlichen Anforderungen entsprechend kenntlich gemacht. Die Versicherung selbstständiger Arbeit gilt auch für Zeichnungen, Skizzen oder grafische Darstellungen. Die Arbeit wurde bisher in gleicher oder ähnlicher Form weder derselben noch einer anderen Prüfungsbehörde vorgelegt und auch noch nicht veröffentlicht. Mit der Abgabe der elektronischen Fassung der endgültigen Version der Arbeit nehme ich zur Kenntnis, dass diese mit Hilfe eines Plagiatserkennungsdienstes auf enthaltene Plagiate überprüft und ausschließlich für Prüfungszwecke gespeichert wird.

\vspace{5cm}

\begin{tabular}{lp{2em}l}
 \hspace{4cm}   && \hspace{5cm} \\\cline{1-1}\cline{3-3}
 Datum     && Unterschrift
\end{tabular}  

\end{comment}

\end{document}

